\chapter{Conclusion}
\label{chap:concl}

\section{Achievements}
This thesis presents a VR memory palace workflow centred on spatial personalisation. The user prepares a layout based on their home, reviews and edits the walls, adds furniture, and uses the resulting environment to place visual mnemonic cues. The development records describe the practical work needed to support this process, including furniture editing, saving and loading, content import, and improvements to menu and wall interactions.

The evaluation compares this personalised workflow with a generic palace within the same participants. It considers recall quantity and order alongside cognitive load, enjoyment, competence, effort, and presence. Pilot development helped identify overly easy materials and an excessive questionnaire burden, leading to revised word lists and a shorter presence measure.

The current numerical snapshot contains a positive mean personalised-minus-generic difference in immediate recall accuracy of 6.28 percentage points across 18 complete pairs. However, the confidence interval crosses zero, the direction of the difference varies across participants, and the workbook remains blocked for interpretation. This draft therefore makes no final claim that personalisation improves recall.

\section{Limitations}
The principal limitation is the unresolved state of the analysis data. Duplicate questionnaire records and unreviewed responses must be reconciled before the results can support a final conclusion. The snapshot also combines materials versions, while the distinction between pilot and main-study sessions requires confirmation.

The experimental contrast combines familiarity with active authoring and extra exposure. It cannot isolate which of those components produces a difference. Follow-up observations are incomplete, actual condition order has not been recovered for the workbook's diagnostic, and nominal follow-up labels do not establish precise retention intervals.

The system description is grounded in the available development logs rather than an audit of the final Unity build. Hardware, software versions, screenshots, recruitment details, and the exact administered questionnaire forms remain necessary additions. These gaps are recorded in Appendix~\ref{app:completion} so that they can be resolved without turning assumptions into facts.

\section{Future Work}
The immediate priority is to complete the existing study record. This includes reviewing response decisions, resolving duplicate records, verifying participant and condition assignments, and identifying the final analysis cohort. The analysis should then be rerun with explicit reporting of missingness, materials versions, and condition order.

A subsequent experiment could separate familiarity from construction. Krokos et al.~\cite{krokos2019} explicitly proposed studying user-built palaces and comparing highly personal with shared palaces; the present study cannot resolve both questions because its home condition also involves authoring and extra exposure. Crossing familiar versus unfamiliar layouts with self-edited versus prepared layouts, while matching exposure time, would estimate those effects separately. Yang et al.'s participant feedback~\cite{yang2021} and Perera et al.'s proposed preference-based templates~\cite{perera2019} also motivate measuring whether users prefer the palace they create. Repeated-use sessions could then examine whether the initial preparation effort becomes worthwhile over several learning tasks and whether reusing loci causes interference.

The broader proposal also includes learning-content personalisation and visual and auditory mnemonic generation. These remain useful directions, but should be evaluated separately from the spatial comparison. Moll and Sykes~\cite{moll2023} allowed a limited choice of images and suggested testing more personally chosen cues. A content study could hold the environment constant while varying cue choice or how the material is segmented and represented. An auditory study could compare visual cues with audio or combined cues while preserving the same target material and exposure time.

Finally, educational use should be evaluated using the outcomes it is intended to improve. Remembering a word list provides a controlled first test, but learning concepts also involves understanding and applying them. A later study should therefore include comprehension and transfer tasks rather than infer educational effectiveness from recall alone.

\section{Closing Remarks}
The project turns the idea of a familiar memory palace into an editable VR activity. Its central question is whether that personal connection to the space improves the experience and the memory task enough to justify the work of creating it. The present draft establishes the system account, evaluation structure, and limits of the available evidence. A final answer to the research question depends on completing and reconciling the study data.

